\section{確率変数と確率分布}
\subsection{授業内容}
\subsubsection{科目の中でのこのコマの位置づけ}
心理統計は個人差や偶然性を表すために確率の考え方が必須である。
しかし確率の概念は数学的に高度かつ抽象的であり，直感的に理解することが難しい。概念的には，空間を標準化したときの相対的な大きさを指標化したもの，すなわち面積の一種と考えるとわかりやすいし，複雑な性質よりも最低限守るべき公理の方からアプローチした方が理解しやすい。ただしこのとき，，確率変数とその実現値の違いに注意する。
とはいえこのコースの目的は，数学的内容を座学で学ぶのではなく，あくまでもプログラミングによる演習を通じて具体的に理解することにある。
本講ではまず概念的な理解の復習からすすめる。
なお以後，\textcite{kosugi2023}を副読本としながら授業を進めていく。
\subsubsection{コマ主題細目}
\begin{description}
	\item[確率の基礎] 高校数学までの範囲で学ぶ確率は，すべての場合の数を全体としたうち，該当する事象の生じる回数を相対頻度で表すものであった。一方我々は，天気予報で降水確率が30\%というように，確率的な表現を日常的にも用いる。天気予報は未来の予測であるから，「起こりうるすべての未来」を書き出して生じる天気の割合を表現できるはずもない。こうした不確実なものであっても，確率という言葉で表現できるようにコルモゴロフが公理として最も基本的な枠組みを提供した。ここでは認識論に深く立ち入ることはせず，確率という数字が従うべき法則について確認する。
		\begin{flushright}
			$\to$ \textcite{Maeda2017}のP.78--83, \textcite{kosugi2023}のP.61--62.
		\end{flushright}
	\item[Rの確率関数] Rには確率関数がいくつか用意されており，接頭語\verb|d,p,q,r|と関数名からなる。ここでは正規分布を例に，\verb|d,p,q-norm| がそれぞれ何を表しているかを確認する。確率は確率密度関数で表される領域の面積であり，積分によって計算されることを近似計算で理解する。
		\begin{flushright}
			$\to$\textcite{kosugi2023}のP.66--77.
		\end{flushright}
	\item[確率変数の期待値と分散] 確率変数の期待値(加重平均)と分散について，定義式を確認する。また正規分布においては期待値が位置パラメータ$\mu$に，分散がスケールパラメータ$\sigma^2$に一致することを確認する。
		\begin{flushright}
			$\to$\textcite{kosugi2023}のP.83--86.
		\end{flushright}
\subsubsection{キーワード}
\begin{itemize}
	\item 確率変数と確率変数の実現値
	\item \verb|dnorm,pnorm,qnorm,rnorm|
	\item 確率分布関数
\end{itemize}
\subsection{授業情報}
\paragraph{コマの展開方法}
講義/遠隔/演習
\subsubsection{予習・復習課題}
\paragraph{予習} 一年時の「データ解析基礎」ですでに確率については学んでいる。テキストや関連資料を見て改めて確率の概念を復習しておくことが，このセクションの予習になる。
\paragraph{復習} Rの正規分布を例に，関数の形とパラメータの関係，期待値と分散について学んできたが，正規分布以外の確率分布関数についても定義や性質を調べつつ，Rの関数をつかって描画・計算して理解を進めておくと良い。